
\documentclass{article}
\usepackage{colortbl}
\usepackage{makecell}
\usepackage{multirow}
\usepackage{supertabular}

\begin{document}

\newcounter{utterance}

\centering \large Interaction Transcript for game `dond', experiment `coop\_en', episode 16 with reflexion{-}Qwen3.5{-}27B.
\vspace{24pt}

{ \footnotesize  \setcounter{utterance}{1}
\setlength{\tabcolsep}{0pt}
\begin{supertabular}{c@{$\;$}|p{.15\linewidth}@{}p{.15\linewidth}p{.15\linewidth}p{.15\linewidth}p{.15\linewidth}p{.15\linewidth}}
    \# & $\;$A & \multicolumn{4}{c}{Game Master} & $\;\:$B\\
    \hline

    \theutterance \stepcounter{utterance}  
    & & \multicolumn{4}{p{0.6\linewidth}}{
        \cellcolor[rgb]{0.9,0.9,0.9}{
            \makecell[{{p{\linewidth}}}]{
                \texttt{\tiny{[P1$\langle$GM]}}
                \texttt{You have attempted this same task before and did not succeed. These are the plans you wrote after those attempts.} \\
\texttt{Trial \#1: New plan: I failed because I assumed the other player would accept my initial offer without verifying their preferences, leading to a situation where they submitted a conflicting secret proposal before I could negotiate. My strategy of immediately claiming all high{-}value items (leaves and magnet) likely left the other player with nothing of value to them, causing a deadlock. Next time, I will first ask the other player what items they value most before making any claims. I will then propose a split that explicitly leaves the items they value for them, ensuring our proposals are complementary. If they submit a secret proposal prematurely, I will immediately counter{-}propose a split that includes the items they likely wanted based on the remaining inventory, rather than just repeating my own greedy claim.} \\
\\ 
\\ 
\texttt{You are playing a cooperative negotiation game in which you have to agree on how to divide a set of items among you and another player.} \\
\\ 
\texttt{Rules:} \\
\texttt{(a) You and the other player are given a set of items. Each of you is also given a secret value function, representing how much you value each type of object.} \\
\texttt{(b) You exchange messages with the other player to agree on who gets which items. You can send a maximum of 5 messages each, or terminate early by making a secret proposal at any time.} \\
\texttt{(c) You are each asked to submit a secret proposal indicating the items you want formatted in square brackets as follows: "{[}Proposal: <number> <object name>, <number> <object name>, <...>{]}"} \\
\texttt{(d) If your proposals are complementary, i.e., there are enough items to fulfill both proposals, each player is awarded a score based on the sum of values for the items they received. Otherwise, both of you get zero points.} \\
\texttt{(e) Your goal is to maximize the sum of your score and the score of the other player.} \\
\\ 
\texttt{Let us start.} \\
\\ 
\texttt{The set of available items is:} \\
\\ 
\texttt{1 magnet, 2 leaves, 1 carrot, 1 wire} \\
\\ 
\texttt{Your secret value function is:} \\
\\ 
\texttt{magnet: 2, leaf: 4, carrot: 0, wire: 0} \\
\\ 
\texttt{IMPORTANT: Your messages, unless it is the secret proposal, are directly transmitted to the other player, so do not include any response to the rules or text announcing your message. To make a secret proposal, use the indicated format. Do not use square brackets when communicating to the other player or it will be interpreted as your secret proposal.} \\
            }
        }
    }
    & \\ \\

    \theutterance \stepcounter{utterance}  
    & \multicolumn{4}{p{0.6\linewidth}}{
        \cellcolor[rgb]{0.9,0.9,0.9}{
            \makecell[{{p{\linewidth}}}]{
                \texttt{\tiny{[P1$\rangle$GM]}}
                \texttt{Hi! To make sure we both get the best possible outcome, could you tell me which items you value the most? I know I really want the leaves and the magnet, but I want to make sure we don't end up with a conflict. What are your preferences for the magnet, leaves, carrot, and wire?} \\
            }
        }
    }
    & & \\ \\

    \theutterance \stepcounter{utterance}  
    & & & \multicolumn{2}{p{0.3\linewidth}}{
        \cellcolor[rgb]{0.9,0.9,0.9}{
            \makecell[{{p{\linewidth}}}]{
                \texttt{\tiny{[GM$|$GM]}}
                \texttt{continue} \\
            }
        }
    }
    & & \\ \\

    \theutterance \stepcounter{utterance}  
    & & & \multicolumn{2}{p{0.3\linewidth}}{
        \cellcolor[rgb]{0.9,0.9,0.9}{
            \makecell[{{p{\linewidth}}}]{
                \texttt{\tiny{[GM$|$GM]}}
                \texttt{Hi! To make sure we both get the best possible outcome, could you tell me which items you value the most? I know I really want the leaves and the magnet, but I want to make sure we don't end up with a conflict. What are your preferences for the magnet, leaves, carrot, and wire?} \\
            }
        }
    }
    & & \\ \\

    \theutterance \stepcounter{utterance}  
    & & \multicolumn{4}{p{0.6\linewidth}}{
        \cellcolor[rgb]{0.9,0.9,0.9}{
            \makecell[{{p{\linewidth}}}]{
                \texttt{\tiny{[GM$\rangle$P2]}}
                \texttt{You have attempted this same task before and did not succeed. These are the plans you wrote after those attempts.} \\
\texttt{Trial \#1: New plan: I failed because I misinterpreted the other player's message as a statement of their own values, when it was actually a description of *my* values (the prompt's secret function) that they were quoting or reacting to. Specifically, the other player stated they value leaves (4) and magnet (2), while I (the AI) actually value carrot (4) and magnet (3). By assuming the other player wanted the leaves and magnet, I proposed a split where they got the items they valued highly, but I gave away the carrot (my highest value) to them, resulting in a low total score.} \\
\\ 
\texttt{Next time, I will:} \\
\texttt{1.  **Strictly separate my secret values from the other player's stated preferences.** I will not assume the other player's message describes their own values unless explicitly stated as "I value...".} \\
\texttt{2.  **Calculate the optimal split based on the actual values:**} \\
\texttt{    *   My values: Carrot (4), Magnet (3), Leaf (1), Wire (1).} \\
\texttt{    *   Other player's stated values: Leaf (4), Magnet (2), Carrot (0), Wire (0).} \\
\texttt{    *   Optimal allocation: I take Carrot (4) and Magnet (3) = 7 points. They take 2 Leaves (8) and Wire (0) = 8 points. Total = 15.} \\
\texttt{3.  **Propose the correct split clearly:** I will explicitly state, "I will take the carrot and magnet, and you take the leaves and wire," before submitting the secret proposal.} \\
\texttt{4.  **Submit the proposal immediately after agreement:** I will ensure the proposal format `{[}Proposal: 1 carrot, 1 magnet{]}` matches the agreed{-}upon items exactly, avoiding any confusion about who gets what.} \\
\\ 
\\ 
\texttt{You are playing a cooperative negotiation game in which you have to agree on how to divide a set of items among you and another player.} \\
\\ 
\texttt{Rules:} \\
\texttt{(a) You and the other player are given a set of items. Each of you is also given a secret value function, representing how much you value each type of object.} \\
\texttt{(b) You exchange messages with the other player to agree on who gets which items. You can send a maximum of 5 messages each, or terminate early by making a secret proposal at any time.} \\
\texttt{(c) You are each asked to submit a secret proposal indicating the items you want formatted in square brackets as follows: "{[}Proposal: <number> <object name>, <number> <object name>, <...>{]}"} \\
\texttt{(d) If your proposals are complementary, i.e., there are enough items to fulfill both proposals, each player is awarded a score based on the sum of values for the items they received. Otherwise, both of you get zero points.} \\
\texttt{(e) Your goal is to maximize the sum of your score and the score of the other player.} \\
\\ 
\texttt{Let us start.} \\
\\ 
\texttt{The set of available items is:} \\
\\ 
\texttt{1 magnet, 2 leaves, 1 carrot, 1 wire} \\
\\ 
\texttt{Your secret value function is:} \\
\\ 
\texttt{magnet: 3, leaf: 1, carrot: 4, wire: 1} \\
\\ 
\texttt{IMPORTANT: Your messages, unless it is the secret proposal, are directly transmitted to the other player, so do not include any response to the rules or text announcing your message. To make a secret proposal, use the indicated format. Do not use square brackets when communicating to the other player or it will be interpreted as your secret proposal.} \\
\\ 
\\ 
\texttt{Hi! To make sure we both get the best possible outcome, could you tell me which items you value the most? I know I really want the leaves and the magnet, but I want to make sure we don't end up with a conflict. What are your preferences for the magnet, leaves, carrot, and wire?} \\
            }
        }
    }
    & \\ \\

    \theutterance \stepcounter{utterance}  
    & & & \multicolumn{4}{p{0.6\linewidth}}{
        \cellcolor[rgb]{0.9,0.9,0.9}{
            \makecell[{{p{\linewidth}}}]{
                \texttt{\tiny{[GM$\langle$P2]}}
                \texttt{Hi! I think there might be a bit of a mix{-}up in your message. You mentioned that *you* really want the leaves and the magnet, but based on the game setup, those are actually the items *I* value the most (Carrot: 4, Magnet: 3, Leaf: 1, Wire: 1).} \\
\\ 
\texttt{To help us find the best split, could you clarify which items **you** value the most? For example, do you value the leaves or the carrot more? Once I know your preferences, I can suggest a division where we both get the items we care about most.} \\
            }
        }
    }
     \\ \\

    \theutterance \stepcounter{utterance}  
    & & & \multicolumn{2}{p{0.3\linewidth}}{
        \cellcolor[rgb]{0.9,0.9,0.9}{
            \makecell[{{p{\linewidth}}}]{
                \texttt{\tiny{[GM$|$GM]}}
                \texttt{continue} \\
            }
        }
    }
    & & \\ \\

    \theutterance \stepcounter{utterance}  
    & & & \multicolumn{2}{p{0.3\linewidth}}{
        \cellcolor[rgb]{0.9,0.9,0.9}{
            \makecell[{{p{\linewidth}}}]{
                \texttt{\tiny{[GM$|$GM]}}
                \texttt{Hi! I think there might be a bit of a mix{-}up in your message. You mentioned that *you* really want the leaves and the magnet, but based on the game setup, those are actually the items *I* value the most (Carrot: 4, Magnet: 3, Leaf: 1, Wire: 1).} \\
\\ 
\texttt{To help us find the best split, could you clarify which items **you** value the most? For example, do you value the leaves or the carrot more? Once I know your preferences, I can suggest a division where we both get the items we care about most.} \\
            }
        }
    }
    & & \\ \\

    \theutterance \stepcounter{utterance}  
    & & \multicolumn{4}{p{0.6\linewidth}}{
        \cellcolor[rgb]{0.9,0.9,0.9}{
            \makecell[{{p{\linewidth}}}]{
                \texttt{\tiny{[P1$\langle$GM]}}
                \texttt{Hi! I think there might be a bit of a mix{-}up in your message. You mentioned that *you* really want the leaves and the magnet, but based on the game setup, those are actually the items *I* value the most (Carrot: 4, Magnet: 3, Leaf: 1, Wire: 1).} \\
\\ 
\texttt{To help us find the best split, could you clarify which items **you** value the most? For example, do you value the leaves or the carrot more? Once I know your preferences, I can suggest a division where we both get the items we care about most.} \\
            }
        }
    }
    & \\ \\

    \theutterance \stepcounter{utterance}  
    & \multicolumn{4}{p{0.6\linewidth}}{
        \cellcolor[rgb]{0.9,0.9,0.9}{
            \makecell[{{p{\linewidth}}}]{
                \texttt{\tiny{[P1$\rangle$GM]}}
                \texttt{Ah, I apologize for the confusion! I must have misread my own value function in my previous thought process. Let me clarify:} \\
\\ 
\texttt{My secret values are:} \\
\texttt{{-} Magnet: 2} \\
\texttt{{-} Leaf: 4} \\
\texttt{{-} Carrot: 0} \\
\texttt{{-} Wire: 0} \\
\\ 
\texttt{So, I actually value the **leaves** the most (4 points), followed by the **magnet** (2 points). I don't care about the carrot or the wire at all (0 points).} \\
\\ 
\texttt{Since you mentioned you value the **carrot** (4) and **magnet** (3) the most, and only 1 point for leaves and wire, here is a proposal that seems perfect for both of us:} \\
\\ 
\texttt{{-} **You** take the **carrot** (your highest value) and the **magnet** (your second highest).} \\
\texttt{{-} **I** take the **2 leaves** (my highest value).} \\
\texttt{{-} We can split the **wire** however you like, or I can take it since it doesn't matter to me, but it doesn't affect our scores.} \\
\\ 
\texttt{This way, you get your top two items, and I get my top item (the leaves). Does this sound good to you? If so, I can submit the secret proposal immediately.} \\
            }
        }
    }
    & & \\ \\

    \theutterance \stepcounter{utterance}  
    & & & \multicolumn{2}{p{0.3\linewidth}}{
        \cellcolor[rgb]{0.9,0.9,0.9}{
            \makecell[{{p{\linewidth}}}]{
                \texttt{\tiny{[GM$|$GM]}}
                \texttt{continue} \\
            }
        }
    }
    & & \\ \\

    \theutterance \stepcounter{utterance}  
    & & & \multicolumn{2}{p{0.3\linewidth}}{
        \cellcolor[rgb]{0.9,0.9,0.9}{
            \makecell[{{p{\linewidth}}}]{
                \texttt{\tiny{[GM$|$GM]}}
                \texttt{Ah, I apologize for the confusion! I must have misread my own value function in my previous thought process. Let me clarify:} \\
\\ 
\texttt{My secret values are:} \\
\texttt{{-} Magnet: 2} \\
\texttt{{-} Leaf: 4} \\
\texttt{{-} Carrot: 0} \\
\texttt{{-} Wire: 0} \\
\\ 
\texttt{So, I actually value the **leaves** the most (4 points), followed by the **magnet** (2 points). I don't care about the carrot or the wire at all (0 points).} \\
\\ 
\texttt{Since you mentioned you value the **carrot** (4) and **magnet** (3) the most, and only 1 point for leaves and wire, here is a proposal that seems perfect for both of us:} \\
\\ 
\texttt{{-} **You** take the **carrot** (your highest value) and the **magnet** (your second highest).} \\
\texttt{{-} **I** take the **2 leaves** (my highest value).} \\
\texttt{{-} We can split the **wire** however you like, or I can take it since it doesn't matter to me, but it doesn't affect our scores.} \\
\\ 
\texttt{This way, you get your top two items, and I get my top item (the leaves). Does this sound good to you? If so, I can submit the secret proposal immediately.} \\
            }
        }
    }
    & & \\ \\

    \theutterance \stepcounter{utterance}  
    & & \multicolumn{4}{p{0.6\linewidth}}{
        \cellcolor[rgb]{0.9,0.9,0.9}{
            \makecell[{{p{\linewidth}}}]{
                \texttt{\tiny{[GM$\rangle$P2]}}
                \texttt{Ah, I apologize for the confusion! I must have misread my own value function in my previous thought process. Let me clarify:} \\
\\ 
\texttt{My secret values are:} \\
\texttt{{-} Magnet: 2} \\
\texttt{{-} Leaf: 4} \\
\texttt{{-} Carrot: 0} \\
\texttt{{-} Wire: 0} \\
\\ 
\texttt{So, I actually value the **leaves** the most (4 points), followed by the **magnet** (2 points). I don't care about the carrot or the wire at all (0 points).} \\
\\ 
\texttt{Since you mentioned you value the **carrot** (4) and **magnet** (3) the most, and only 1 point for leaves and wire, here is a proposal that seems perfect for both of us:} \\
\\ 
\texttt{{-} **You** take the **carrot** (your highest value) and the **magnet** (your second highest).} \\
\texttt{{-} **I** take the **2 leaves** (my highest value).} \\
\texttt{{-} We can split the **wire** however you like, or I can take it since it doesn't matter to me, but it doesn't affect our scores.} \\
\\ 
\texttt{This way, you get your top two items, and I get my top item (the leaves). Does this sound good to you? If so, I can submit the secret proposal immediately.} \\
            }
        }
    }
    & \\ \\

    \theutterance \stepcounter{utterance}  
    & & & \multicolumn{4}{p{0.6\linewidth}}{
        \cellcolor[rgb]{0.9,0.9,0.9}{
            \makecell[{{p{\linewidth}}}]{
                \texttt{\tiny{[GM$\langle$P2]}}
                \texttt{That sounds like a perfect plan! Let's confirm the split to ensure we both get the maximum points:} \\
\\ 
\texttt{*   **You:** 1 Carrot (4 pts) + 1 Magnet (3 pts) = **7 points**} \\
\texttt{*   **Me:** 2 Leaves (4 pts each) = **8 points**} \\
\texttt{*   **Wire:** Since it's worth 0 to both of us, I'll take it just to clear the list, but it doesn't change the score.} \\
\\ 
\texttt{This gives us a combined total of 15 points, which is the best possible outcome.} \\
\\ 
\texttt{I agree to this split. I will now submit my secret proposal.} \\
\\ 
\texttt{{[}Proposal: 1 carrot, 1 magnet, 1 wire{]}} \\
            }
        }
    }
     \\ \\

    \theutterance \stepcounter{utterance}  
    & & & \multicolumn{2}{p{0.3\linewidth}}{
        \cellcolor[rgb]{0.9,0.9,0.9}{
            \makecell[{{p{\linewidth}}}]{
                \texttt{\tiny{[GM$|$GM]}}
                \texttt{proposal} \\
            }
        }
    }
    & & \\ \\

    \theutterance \stepcounter{utterance}  
    & & & \multicolumn{2}{p{0.3\linewidth}}{
        \cellcolor[rgb]{0.9,0.9,0.9}{
            \makecell[{{p{\linewidth}}}]{
                \texttt{\tiny{[GM$|$GM]}}
                [1, 0, 1, 1]
            }
        }
    }
    & & \\ \\

    \theutterance \stepcounter{utterance}  
    & & \multicolumn{4}{p{0.6\linewidth}}{
        \cellcolor[rgb]{0.9,0.9,0.9}{
            \makecell[{{p{\linewidth}}}]{
                \texttt{\tiny{[P1$\langle$GM]}}
                \texttt{The other player has submitted a secret proposal. It is now your turn to submit a proposal of your own.} \\
            }
        }
    }
    & \\ \\

    \theutterance \stepcounter{utterance}  
    & \multicolumn{4}{p{0.6\linewidth}}{
        \cellcolor[rgb]{0.9,0.9,0.9}{
            \makecell[{{p{\linewidth}}}]{
                \texttt{\tiny{[P1$\rangle$GM]}}
                \texttt{{[}Proposal: 2 leaf, 1 wire{]}} \\
            }
        }
    }
    & & \\ \\

    \theutterance \stepcounter{utterance}  
    & & & \multicolumn{2}{p{0.3\linewidth}}{
        \cellcolor[rgb]{0.9,0.9,0.9}{
            \makecell[{{p{\linewidth}}}]{
                \texttt{\tiny{[GM$|$GM]}}
                \texttt{proposal} \\
            }
        }
    }
    & & \\ \\

    \theutterance \stepcounter{utterance}  
    & & & \multicolumn{2}{p{0.3\linewidth}}{
        \cellcolor[rgb]{0.9,0.9,0.9}{
            \makecell[{{p{\linewidth}}}]{
                \texttt{\tiny{[GM$|$GM]}}
                [0, 2, 0, 1]
            }
        }
    }
    & & \\ \\

    \theutterance \stepcounter{utterance}  
    & & & \multicolumn{2}{p{0.3\linewidth}}{
        \cellcolor[rgb]{0.9,0.9,0.9}{
            \makecell[{{p{\linewidth}}}]{
                \texttt{\tiny{[GM$|$GM]}}
                [[0, 2, 0, 1], [1, 0, 1, 1]]
            }
        }
    }
    & & \\ \\

    \theutterance \stepcounter{utterance}  
    & \multicolumn{4}{p{0.6\linewidth}}{
        \cellcolor[rgb]{0.95,0.95,0.95}{
            \makecell[{{p{\linewidth}}}]{
                \texttt{\tiny{[P2$\rightarrow$P1]}}
                \texttt{Post{-}episode Reflexion (not sent to GM):} \\
\texttt{New plan: I failed because I misinterpreted the other player's feedback as a simple clarification rather than a critical correction of my initial strategy. The other player explicitly stated they value the **carrot** and **magnet** most, yet I proceeded to propose a split where I took the **leaves** and **wire**, implicitly leaving the magnet for them but failing to explicitly confirm the magnet's allocation in a way that guaranteed complementarity before they acted. More critically, I did not wait for their explicit agreement to the specific split before they submitted their own secret proposal, which likely conflicted with my assumption. Next time, after the other player reveals their values, I will explicitly state the exact division (e.g., "You take Carrot and Magnet; I take Leaves") and ask for a direct confirmation ("Does this work?") before submitting any secret proposal. If they submit a proposal first, I will immediately analyze their specific item list and submit a proposal that takes *only* the remaining items, ensuring no overlap, rather than assuming a split based on my own preferences.} \\
            }
        }
    }
    & & \\ \\

    \theutterance \stepcounter{utterance}  
    & & & \multicolumn{4}{p{0.6\linewidth}}{
        \cellcolor[rgb]{0.95,0.95,0.95}{
            \makecell[{{p{\linewidth}}}]{
                \texttt{\tiny{[P1$\rightarrow$P2]}}
                \texttt{Post{-}episode Reflexion (not sent to GM):} \\
\texttt{New plan: I failed because I confused the roles in the final proposal. The other player correctly identified that I (the AI) should take the carrot and magnet, and they should take the leaves. However, in my secret proposal, I listed the items I *wanted* (carrot, magnet, wire) but failed to explicitly assign the leaves to the other player in a way that ensured a complementary split, or I misinterpreted the "complementary" rule by not ensuring the other player's proposal would cover the remaining items. More critically, looking at the dialogue, the other player proposed: "You take carrot and magnet, I take 2 leaves." I agreed, but my proposal `{[}Proposal: 1 carrot, 1 magnet, 1 wire{]}` left the 2 leaves unassigned in *my* proposal. If the other player's proposal was `{[}Proposal: 2 leaves{]}`, the split is valid. However, the error likely stems from the fact that I included the "wire" in my proposal when the other player said "I can take it" or "split it however you like," potentially causing a conflict if they also claimed it, or simply that I didn't explicitly confirm the *other* player's claim to the leaves in the final step.} \\
\\ 
\texttt{Actually, re{-}reading the failure: The other player said "I take the 2 leaves". I said "I'll take the wire". My proposal was `{[}Proposal: 1 carrot, 1 magnet, 1 wire{]}`. The other player likely proposed `{[}Proposal: 2 leaves{]}`. This *should* have been complementary (Carrot, Magnet, Wire vs Leaves). Why did it fail?} \\
\texttt{Ah, the other player said "I take the 2 leaves... We can split the wire however you like, or I can take it". If the other player decided to take the wire in their secret proposal (since they said "I can take it"), and I also took the wire in my proposal, we both claimed the wire. Since there is only 1 wire, the proposals are not complementary (overlap on wire), resulting in 0 points.} \\
\\ 
\texttt{**Corrected Strategy:**} \\
\texttt{1.  **Resolve Ambiguity on Low{-}Value Items:** When the other player says "I can take it" regarding an item (like the wire), I must explicitly assign it to *them* in my agreement to avoid overlap. I should say, "Okay, you take the wire too," rather than "I'll take it."} \\
\texttt{2.  **Verify Complementary Sets:** Before submitting, I must mentally check: "If I take X, Y, Z, does the other player's stated plan take the *exact* remaining items?"} \\
\texttt{3.  **Final Proposal Construction:**} \\
\texttt{    *   My items: Carrot, Magnet.} \\
\texttt{    *   Their items: 2 Leaves, Wire (since they offered to take it).} \\
\texttt{    *} \\
            }
        }
    }
     \\ \\

\end{supertabular}
}

\end{document}
